\section{Strong resonance currents and first-order transport of all five moments}
\label{cur:section}

The preceding section defined the transport coefficients through a weighted
cell problem. We now identify them for the solution family of
Theorem \ref{hyd:main}. The key quantity is the microscopic entropy
isolated in Proposition \ref{hyd:entropy-geometry}. Its exact time balance
relates collision dissipation to the residual flux and determines the
norm of the resonance current after its weak limit has been found.
We retain the same initial data and prescribed Euler interval throughout.
To fix the limiting measure, this section uses the local abbreviation
$N=\bar N=N_{\bar\theta}$; the actual moment-matched equilibrium is
always denoted by $N_c=N_{\theta_c}$. Set
\begin{equation}\label{cur:variables}
 d_c=\frac{f_c}{N_c},\qquad w_c=1-d_c^{-1},\qquad
 z_c=cw_c,\qquad F_E=-\mathcal T\log N.
\end{equation}
The variable $w_c$ measures the relative deviation of the reciprocal
distribution, since $1/f_c=(1-w_c)/N_c$. It gives
$\Delta(1/f_c)=-\Delta(w_c/N_c)$, while the factor $c$ extracts the
first-order microscopic response.
By \eqref{ons:euler-force}, $F_E\in Q_NL^\infty(D)$.
Use the weighted inverse of Lemma \ref{ons:cell} to define
\begin{equation}\label{cur:limitcell}
 L_Nu=F_E,\qquad \int_D N\Psi u\dd k=0.
\end{equation}
This cell problem depends on $(t,X)$, and all five normalization
integrals are well defined. The tensor $\mathsf K$ below has the
$c$ normalization in \eqref{ons:definition}. Write
\[
 \mathcal H=L^2((0,T)\times\T^3;H_\nu),\qquad
 \mathcal B=L^2((0,T)\times\T^3;L^1(D)).
\]

\begin{theorem}[Strong limits of resonance currents and moment fluxes]\label{cur:main}
Under the assumptions of Theorem \ref{hyd:main}, the following limits
hold on the entire prescribed interval $(0,T)$:
\begin{align}
 z_c&\longrightarrow u
       &&\text{strongly in }\mathcal H,
       \label{cur:microreciprocal}\\
 c(f_c-N_c)&\longrightarrow Nu
       &&\text{strongly in }\mathcal H\cap\mathcal B.
       \label{cur:microphysical}
\end{align}
Define the signed four-leg resonance current
\[
 \mathcal J_c=\frac c2\sqrt{\prod_{i=0}^3f_{c,i}}\,
                    \Delta(1/f_c).
\]
Then
\begin{equation}\label{cur:resonanceconclusion}
 \mathcal J_c\longrightarrow
 -\frac12\sqrt{\prod_{i=0}^3N_i}\,\Delta(u/N)
 \quad\text{strongly in }L^2(\dd t\dd X\dd\Xi_D),
\end{equation}
and the physical collision term satisfies
\begin{equation}\label{cur:collisionconclusion}
 c\mathcal C(f_c)\longrightarrow\mathcal TN
       \quad\text{in }L^2((0,T)\times\T^3\times D).
\end{equation}
For $i=1,2,3$ and $a=0,\ldots,4$, set
\[
 Z_{c,i}^a=\int_Dv_i\Psi_a(f_c-N_c)\dd k.
\]
All fifteen components satisfy
\begin{equation}\label{cur:fluxconclusion}
 cZ_{c,i}^a\longrightarrow
       \sum_{j=1}^3\sum_{b=0}^4
          \mathsf K_{ij}^{ab}(N)\partial_j\bar\theta_b
       \quad\text{strongly in }L^2((0,T)\times\T^3).
\end{equation}
Consequently, the actual five-moment parameter obeys the first-order law
\begin{equation}\label{cur:firstorder}
 \begin{split}
 &M(\theta_c)\partial_t\theta_c
       +\sum_iJ_i(\theta_c)\partial_i\theta_c
 \\
 &\qquad=\frac1c\sum_{i,j}\partial_i
       \bigl[\mathsf K_{ij}(N_c)\partial_j\theta_c\bigr]
       +o(c^{-1}).
 \end{split}
\end{equation}
The remainder belongs to $L^2(0,T;H^{-1}(\T^3;\R^5))$, with norm
$o(c^{-1})$.
\end{theorem}

The proof has five steps. Theorem \ref{hyd:main} first gives weighted
microscopic bounds. The kinetic equation then identifies the weak limit
on a fixed resonance measure. The entropy balance determines the current
norm, so weak convergence becomes strong convergence. An exact Gram
equation from moment matching recovers the five-dimensional null
component left undetermined by dissipation. Finally, the adjoint of the
collision-difference operator returns the physical collision term,
and integration gives all five moment fluxes. Ratios in the corner
layer need only fixed positive upper and lower bounds.

\subsection{Actual moment matching and the entropy identity}

Theorem \ref{hyd:main} and the decomposition $f_c-N_c=e_c+G_c$ in its
proof give positive constants $d_-,d_+$, independent of $c$, such that
\begin{align}
 d_-\le d_c\le d_+,&\qquad
 \|\theta_c-\bar\theta\|_{L^\infty_tH^2_X}\le Cc^{-1},
          \label{cur:inputmacro}\\
 \|f_c-N_c\|_{L^\infty_tL^2_{X,k}}&\le Cc^{-3/4},\qquad
 \|f_c-N_c\|_{L^2_{t,X}H_\nu}\le Cc^{-1}.
          \label{cur:inputmicro}
\end{align}
The three-dimensional embedding $H^2_X\hookrightarrow C^0_X$ yields
uniform convergence $N_c/N\to1$, and
$\nabla_X\theta_c\to\nabla_X\bar\theta$ strongly in $L^2_{t,X}$.
All parameters remain in a common compact subset of the positive domain.

Define
\begin{align}
 H_c(t)&=\int_{\T^3}\mathscr H(f_c\mid N_c)\dd X
       =\int_{\T^3\times D}\Phi(d_c)\dd X\dd k,
       \qquad\Phi(d)=d-1-\log d,\label{cur:entropydef}\\
 D_c(t)&=\frac14\int_{\T^3}\int
       \prod_i f_{c,i}\,[\Delta(1/f_c)]^2\dd\Xi_D\dd X.
       \label{cur:dissdef}
\end{align}
The ratio bounds and \eqref{cur:inputmicro} imply
\begin{equation}\label{cur:entropybound}
 \sup_{t\le T}H_c(t)\le Cc^{-3/2},\qquad
 \|z_c\|_{L^2_{t,X}H_\nu}\le C,\qquad
 \sup_{t\le T}c\int_{X,k}|w_c|^2\le Cc^{-1/2}.
\end{equation}
Applying \eqref{hyd:entropy-pythagoras} to the actual moment match gives
\begin{equation}\label{cur:entropy-splitting}
 \int_{\T^3}\mathscr H(f_c\mid N)\dd X
       =H_c(t)+\int_{\T^3}\mathscr H(N_c\mid N)\dd X.
\end{equation}
By \eqref{hyd:entropy-comparison} and \eqref{cur:inputmacro}, the second
term is uniformly $O(c^{-2})$. We use the balance for $H_c$ itself
to retain the exact relation between microscopic dissipation and
the actual moment fluxes.
In particular, $w_c$ is uniformly bounded and tends to zero in volume
measure. The difference between the two microscopic coordinates satisfies
\begin{equation}\label{cur:coordinatedifference}
 c[(d_c-1)-w_c]=cd_cw_c^2\longrightarrow0
                 \quad\text{in }L^\infty_tL^1_{X,k}.
\end{equation}

All entropy balances below are integrated from the given initial time.
Moment matching gives $\int_D\Psi(f_c-N_c)\dd k=0$, and hence
\[
 H_c=\int_{X,k}-\log f_c-\int_{X,k}-\log N_c.
\]
The four-leg symmetry of \eqref{geom:collision} and periodic transport
give $-cD_c$ for the time derivative of the first term.
For the second, set
$\mathcal E(U(\theta))=\int_D-\log N_\theta\dd k$.
Equation \eqref{hyd:entropy-hessians} gives
$\nabla_U\mathcal E=-\theta$. The moment equation for the same solution is
\begin{equation}\label{cur:exactmoments}
 \partial_tU(\theta_c)+\sum_i\partial_i\mathcal F_i(\theta_c)
                  =-\sum_i\partial_iZ_{c,i}.
\end{equation}
The equilibrium entropy flux is $\int_D-v_i\log N_\theta\dd k$,
whose divergence integrates to zero on the torus. Therefore
\[
 \frac{\dd}{\dd t}\int_X\mathcal E(U(\theta_c))
             =-\sum_i\int_X\partial_i\theta_c\cdot Z_{c,i}.
\]
Combining the two derivatives gives
\begin{equation}\label{cur:entropyidentity}
 H_c'+cD_c=\sum_i\int_X\partial_i\theta_c\cdot Z_{c,i},
              \qquad H_c(0)=0.
\end{equation}
The positive classical solution of Theorem \ref{hyd:main} justifies
every differentiation. Passing this identity to the limit will use
only the parameter and spatial-gradient convergence in
\eqref{cur:inputmacro}; time derivatives of the actual parameters
have already been eliminated by the exact calculation.

\begin{lemma}[Strong convergence of bounded multipliers]\label{cur:multiplierlemma}
 Let $(\Omega,\lambda)$ be a finite measure space. Suppose $b_n\to b$
 in $\lambda$-measure and $|b_n|,|b|\le C$.
 If $\sigma\ll\lambda$ is a finite positive measure, then for every
 fixed $V\in L^2(\sigma)$,
 \[
                 \|(b_n-b)V\|_{L^2(\sigma)}\longrightarrow0.
 \]
\end{lemma}
\begin{proof}
 For $\delta>0$, let $A_n=\{|b_n-b|>\delta\}$.
 The finite measure $|V|^2\sigma$ is absolutely continuous with respect
 to $\lambda$, and $\lambda(A_n)\to0$. Thus
 \[
 \|(b_n-b)V\|_{L^2(\sigma)}^2
 \le\delta^2\|V\|_{L^2(\sigma)}^2
                  +4C^2\int_{A_n}|V|^2\dd\sigma.
 \]
 Let $n\to\infty$ and then $\delta\downarrow0$.
\end{proof}

\subsection{A common resonance measure and the limiting weak equation}

On the Hilbert space $\mathcal H$ defined above, fix the measure and operator
\begin{equation}\label{cur:jointmeasure}
 \dd\mu=\prod_{i=0}^3N_i\,\dd t\dd X\dd\Xi_D,
      \qquad \mathsf R_Nh=\tfrac12\Delta(h/N).
\end{equation}
Here $N$ is the fixed limiting distribution, which still varies with
$(t,X)$. The common measure places every resonance current in the
same Hilbert space. Changes in the actual distributions will appear
as explicit multipliers. Four-leg symmetry gives the exact marginal
identity
\begin{equation}\label{cur:marginal}
 \int|b(t,X,k_i)|^2\dd\mu
     =\int_{t,X,k}N^2\nu_N|b|^2\dd t\dd X\dd k,
                   \qquad i=0,1,2,3.
\end{equation}
Consequently $\mathsf R_N:\mathcal H\to L^2(\mu)$ is bounded.
The identity also shows that changing a momentum function on a volume
null set does not alter the $\mu$-equivalence class of any leg.
Thus the operator is defined on the whole weighted Hilbert space.

Set
\begin{equation}\label{cur:multipliers}
 y_c=\frac N{N_c}z_c,\qquad
 m_c=\prod_{i=0}^3\frac{f_{c,i}}{N_i}.
\end{equation}
The variable $y_c$ replaces the matched denominator by the reference
denominator, so $y_c/N=z_c/N_c$. The multiplier $m_c$ records the full
change of the four-leg product relative to the reference measure.
Each $f_c/N-1$ tends to zero in volume $L^2$ and is uniformly bounded.
Applying \eqref{cur:marginal} and subtracting the four factors in turn gives
\begin{equation}\label{cur:coefficientconvergence}
 m_c\to1,\quad \sqrt{m_c}\to1\quad\text{in }L^2(\mu),
          \qquad 0<m_-\le m_c\le m_+<\infty.
\end{equation}
All four legs come from one distribution at the same $(t,X)$.
Equation \eqref{cur:marginal} transfers volume convergence to the
resonance measure. Moreover, for every fixed $V\in L^2(\mu)$,
\begin{equation}\label{cur:fixedmultiplier}
 (m_c-1)V\to0,\qquad(\sqrt{m_c}-1)V\to0
                         \quad\text{in }L^2(\mu).
\end{equation}
This is Lemma \ref{cur:multiplierlemma} with $\lambda=\sigma=\mu$.
It places the multiplier on the fixed test function, allowing the
microscopic vector on the other side to converge only weakly.

By \eqref{cur:entropybound}, extract a subsequence
$z_c\rightharpoonup u_0$ in $\mathcal H$.
Uniform macroscopic convergence also gives $y_c\rightharpoonup u_0$.
Since $\nu_*^{-1}\in L^1(D)$, there is a continuous embedding
\begin{equation}\label{cur:weighteddual}
 \mathcal H\hookrightarrow\mathcal B.
\end{equation}
Apply \eqref{cur:coordinatedifference} to the identity
$\int N_c\Psi d_cz_c\dd k=0$ and let $c\to\infty$. This yields
\begin{equation}\label{cur:limitmoments}
 \int_D N\Psi u_0\dd k=0\quad\text{for almost every }(t,X).
\end{equation}
More explicitly, first pair with an arbitrary smooth space-time function.
Then use uniform convergence $N_c\to N$ and the admissible weighted
dual test $N\Psi/\sqrt{\nu_*}$.

Take a real test function $\psi(t,X,k)=\chi(t,X)b(k)$, where $\chi$
is smooth with time support in $(0,T)$ and $b$ is real and bounded.
Test the kinetic equation with the fixed function $\psi/N$.
Since $\Delta(1/N_c)=0$, its four-leg weak form gives exactly
\begin{equation}\label{cur:weakidentity}
 \int c\mathcal C(f_c)\frac\psi N\dd t\dd X\dd k
           =-\int m_c\mathsf R_Ny_c\,\mathsf R_N\psi\dd\mu.
\end{equation}
On the left, substitute $\mathcal Tf_c=c\mathcal C(f_c)$ and integrate
by parts. Strong volume $L^2$ convergence of $f_c$ gives the limit
$\int(\mathcal TN)\psi/N$; only $\chi$ and the fixed $N$ are differentiated.
By \eqref{cur:fixedmultiplier} and weak convergence, the right-hand side
tends to $-\int\mathsf R_Nu_0\mathsf R_N\psi\dd\mu$.
Using a countable dense family of bounded functions $b$ gives the
almost-everywhere variational equation
\[
 \mathfrak a_N(u_0,b)=\int_DF_Eb\dd k.
\]
As $F_E/\sqrt{\nu_*}\in L^2$, continuity of the form extends the equation
to every real test in $H_\nu$, and then to complex tests by the
sesquilinear convention. Equation \eqref{cur:limitmoments} and weighted
coercivity imply $u_0=u$, with $u$ defined by \eqref{cur:limitcell}.
Every weakly convergent subsequence therefore has the same limit.

\subsection{Norm convergence and recovery of the five-dimensional kernel}

By \eqref{cur:dissdef} and
$\Delta(w_c/N_c)=-\Delta(1/f_c)$,
\begin{equation}\label{cur:normidentity}
 c^2\int_0^TD_c(t)\dd t
             =\|\sqrt{m_c}\mathsf R_Ny_c\|_{L^2(\mu)}^2.
\end{equation}
Multiply \eqref{cur:entropyidentity} by $c$ and integrate from the
initial time to $T$:
\begin{equation}\label{cur:entropynorm}
 cH_c(T)+c^2\int_0^TD_c
        =\int_{t,X,k}z_cF_c,\qquad
 F_c=N_cd_c(v\cdot\nabla_Xp_c),\quad p_c=N_c^{-1}=\theta_c\cdot\Psi.
\end{equation}
The first term tends to zero by \eqref{cur:entropybound}.
Set $p=N^{-1}=\bar\theta\cdot\Psi$ and $F=N(v\cdot\nabla_Xp)$.
Then
\begin{align}
 F_c-F={}&N_cd_cv\cdot\nabla_X(p_c-p)
       +(N_c-N)d_c(v\cdot\nabla_Xp)\notag\\
       &+N(d_c-1)(v\cdot\nabla_Xp).
       \label{cur:workdifference}
\end{align}
After division by $\sqrt{\nu_*}$, the first term is controlled by
$\|\nabla_X\theta_c-\nabla_X\bar\theta\|_{L^2_{t,X}}$ and the second
by uniform macroscopic convergence. In the third term, $d_c-1$ is
uniformly bounded and tends to zero in volume measure, while the
fixed function $|N(v\cdot\nabla_Xp)|^2/\nu_*$ is integrable.
Lemma \ref{cur:multiplierlemma}, applied with volume measure and this
density, gives
\begin{equation}\label{cur:workstrong}
 F_c/\sqrt{\nu_*}\longrightarrow F/\sqrt{\nu_*}
                  \quad\text{strongly in }L^2_{t,X,k}.
\end{equation}
The right-hand side of \eqref{cur:entropynorm} therefore tends to
$\int uF$. The five moment normalizations allow us to add
$N\partial_tp$, so
\[
 \int uF=\int uF_E
     =\int_0^T\!\int_X\mathfrak a_N(u,u)
     =\|\mathsf R_Nu\|_{L^2(\mu)}^2.
\]
On the other hand, \eqref{cur:fixedmultiplier} gives
$\sqrt{m_c}\mathsf R_Ny_c\rightharpoonup\mathsf R_Nu$.
Equations \eqref{cur:normidentity}--\eqref{cur:entropynorm}
give convergence of the norms. Hence
\begin{equation}\label{cur:strongresonance}
 \sqrt{m_c}\mathsf R_Ny_c\longrightarrow\mathsf R_Nu
                         \quad\text{strongly in }L^2(\mu).
\end{equation}
Removing the density in \eqref{cur:jointmeasure} gives
\eqref{cur:resonanceconclusion}. Its sign follows from
$\Delta(w_c/N_c)=-\Delta(1/f_c)$.

Strong convergence of the resonance difference controls only the
complement of the null space: adding
$N\Psi\cdot\zeta^{\mathrm{ker}}(t,X)$ to $y_c$ leaves the difference
unchanged. The null component must therefore be recovered from actual
moment matching. The positive lower bound on $m_c$ and
\eqref{cur:fixedmultiplier}, applied to the fixed vector $\mathsf R_Nu$,
first yield $\mathsf R_N(y_c-u)\to0$ strongly in $L^2(\mu)$.
Write $\kappa=N\Psi$ and $M_N=\int_D\kappa\kappa^T\dd k$.
The extension of the unweighted projection is
\[
 P_Nh=\kappa^TM_N^{-1}\int_D\kappa h\dd k.
\]
It is bounded on $H_\nu$ because $\nu_*^{-1}\in L^1$.
Decompose
\[
 y_c-u=g_c+\kappa^T\zeta_c^{\mathrm{ker}},\qquad g_c=Q_N(y_c-u).
\]
The identity $\mathfrak a_N(Q_Nh,Q_Nh)=\mathfrak a_N(h,h)$ and
\eqref{geom:microcoercivity} imply $g_c\to0$ strongly in $\mathcal H$.
The projection here is the extension of the unweighted one.

Actual moment matching gives
\begin{equation}\label{cur:unknowngram}
 \int_D\kappa b_cy_c\dd k=0,\qquad
 b_c=d_c(N_c/N)^2,
\end{equation}
and consequently
\begin{equation}\label{cur:gramrecovery}
 \left(\int_Db_c\kappa\kappa^T\dd k\right)\zeta_c^{\mathrm{ker}}
      =-\int_D\kappa b_cg_c\dd k
       -\int_D\kappa(b_c-1)u\dd k.
\end{equation}
The matrix on the left is uniformly positive, since
$\int b_c\kappa\kappa^T\ge(\inf b_c)M_N$.
The uniform $L^2$ bound on $\kappa/\sqrt{\nu_*}$ gives, at each $(t,X)$,
\begin{equation}\label{cur:gramestimate}
 |\zeta_c^{\mathrm{ker}}|\le C\bigl(\|g_c\|_{H_\nu}
                         +\|(b_c-1)u\|_{H_\nu}\bigr).
\end{equation}
Since $b_c$ is uniformly bounded and tends to one in volume measure,
Lemma \ref{cur:multiplierlemma} gives
$\|(b_c-1)u\|_{\mathcal H}\to0$.
Thus $\zeta_c^{\mathrm{ker}}\to0$ in $L^2_{t,X}$ and
$y_c\to u$ strongly in $\mathcal H$. Finally,
\[
 z_c-u=(N_c/N)(y_c-u)+(N_c/N-1)u
\]
proves \eqref{cur:microreciprocal}. Decomposing
$c(f_c-N_c)=N_cd_cz_c$ into a strongly convergent difference and a
bounded multiplier acting on the fixed $u$ gives strong convergence
in $\mathcal H$. The embedding \eqref{cur:weighteddual} then gives
\eqref{cur:microphysical}.

\subsection{The physical collision term and the first-order macroscopic equation}

On the physical space $L^2_{t,X,k}$, define
$S_N\varphi=\frac12\Delta\varphi$ with values in $L^2(\mu)$.
Equation \eqref{cur:marginal} gives
\begin{equation}\label{cur:returnbound}
 \|S_N\varphi\|_{L^2(\mu)}^2
       \le4\int_{t,X,k}N^2\nu_N|\varphi|^2
       \le C\|\varphi\|_{L^2_{t,X,k}}^2.
\end{equation}
Its Hilbert adjoint $S_N^*:L^2(\mu)\to L^2_{t,X,k}$ is therefore bounded.
The four-leg weak identity is
\begin{equation}\label{cur:returnidentity}
 c\mathcal C(f_c)=-S_N^*[m_c\mathsf R_Ny_c].
\end{equation}
Equation \eqref{cur:strongresonance} and multiplier convergence on the
fixed limit give $m_c\mathsf R_Ny_c\to\mathsf R_Nu$ strongly in $L^2(\mu)$.
Hence
\[
 c\mathcal C(f_c)\longrightarrow-S_N^*\mathsf R_Nu
                    =-NL_Nu=\mathcal TN,
\]
which proves \eqref{cur:collisionconclusion}.

Every $v_i\Psi_a$ is bounded on $D$, so
\eqref{cur:microphysical} gives the strong limits of all five moment
fluxes simultaneously. By \eqref{ons:cell-flux},
\[
 \int_Dv_i\Psi_aNu\dd k
       =\sum_{j,b}\mathsf K_{ij}^{ab}(N)\partial_j\bar\theta_b,
\]
proving \eqref{cur:fluxconclusion}.
Smoothness of the tensor on compact positive parameter sets and
\eqref{cur:inputmacro} also give
\begin{equation}\label{cur:actualparameterflux}
 \mathsf K(N_c)\nabla_X\theta_c
       -\mathsf K(N)\nabla_X\bar\theta\longrightarrow0
                       \quad\text{strongly in }L^2_{t,X}.
\end{equation}
The bound is the uniform parameter difference times the norm of the
fixed $\nabla_X\bar\theta$, plus
$C\|\nabla_X\theta_c-\nabla_X\bar\theta\|_2$.
Write the flux difference as
\[
 Z_{c,i}=c^{-1}\sum_j\mathsf K_{ij}(N_c)\partial_j\theta_c
                  +\mathfrak R_{c,i}^{\mathrm{flux}},\qquad
 \|\mathfrak R_c^{\mathrm{flux}}\|_{L^2_{t,X}}=o(c^{-1}).
\]
Denote its spatial divergence by
$\mathfrak R_c^{\mathrm{mac}}=\sum_i\partial_i\mathfrak R_{c,i}^{\mathrm{flux}}$.
Then $\|\mathfrak R_c^{\mathrm{mac}}\|_{L^2_tH^{-1}_X}=o(c^{-1})$.
Substitute into \eqref{cur:exactmoments} and use
$DU=-M$, $D\mathcal F_i=-J_i$ to obtain \eqref{cur:firstorder};
spatial divergence maps $L^2_X$ continuously into $H^{-1}_X$.
In conserved variables, the same law reads
\[
 \partial_tU(\theta_c)+\sum_i\partial_i\mathcal F_i(\theta_c)
 =-c^{-1}\sum_{i,j}\partial_i
      [\mathsf K_{ij}(N_c)\partial_j\theta_c]
      -\mathfrak R_c^{\mathrm{mac}}\quad\text{in }L^2_tH^{-1}_X.
\]
The mass row also satisfies
$Z_{c,i}^0=2\int k^{(i)}(f_c-N_c)\dd k=0$, in agreement with the zero
mass row of the tensor. This completes the proof of
Theorem \ref{cur:main}.

\begin{corollary}[Macroscopic entropy and first-order dissipation]\label{cur:entropyproduction}
 Define the equilibrium entropy and entropy flux by
 \[
 \mathcal E(U(\theta))=-\int_D\log N_\theta\dd k,\qquad
 \mathcal Q_i(\theta)=-\int_Dv_i\log N_\theta\dd k.
 \]
 Under the assumptions of Theorem \ref{cur:main},
 \begin{equation}\label{cur:entropyproductionlimit}
 c^2D_c(t)\longrightarrow
 \int_{\T^3}\sum_{i,j}
 (\partial_i\bar\theta)^{\mathsf T}
       \mathsf K_{ij}(N)\partial_j\bar\theta\dd X
       \quad\text{in }L^1(0,T).
 \end{equation}
 The actual parameter also satisfies the local entropy law
 \begin{equation}\label{cur:macroentropy}
 \begin{split}
 &\partial_t\mathcal E(U(\theta_c))
   +\sum_i\partial_i\bigl[\mathcal Q_i(\theta_c)
                                      -\theta_c\cdot Z_{c,i}\bigr]\\
 &\qquad=-\frac1c\sum_{i,j}(\partial_i\theta_c)^{\mathsf T}
                  \mathsf K_{ij}(N_c)\partial_j\theta_c+r_c^{\mathcal E},
 \qquad\|r_c^{\mathcal E}\|_{L^1_{t,X}}=o(c^{-1}).
 \end{split}
 \end{equation}
\end{corollary}
\begin{proof}
 Set $J_c^\mu=\sqrt{m_c}\mathsf R_Ny_c$ and $J^\mu=\mathsf R_Nu$.
 Strong convergence of the resonance current and Cauchy's inequality give
 \[
 \bigl\||J_c^\mu|^2-|J^\mu|^2\bigr\|_{L^1(\mu)}
 \le(\|J_c^\mu\|_2+\|J^\mu\|_2)\|J_c^\mu-J^\mu\|_2\longrightarrow0.
 \]
 Integrate over $(X,k_0,k_1,k_2,k_3)$ and use the fixed-time version of
 \eqref{cur:normidentity}, together with
 \eqref{ons:quadratic-form} and \eqref{ons:euler-force}, to obtain
 \eqref{cur:entropyproductionlimit}.

 Since $\nabla_U\mathcal E=-\theta$ and
 $D_\theta\mathcal Q_i=\mathcal F_i^{\mathsf T}$, the equilibrium
 flux satisfies
 $\theta\cdot\partial_i\mathcal F_i=-\partial_i\mathcal Q_i$.
 Multiply \eqref{cur:exactmoments} by $-\theta_c$ and apply the chain rule:
 \[
 \partial_t\mathcal E(U(\theta_c))
 +\sum_i\partial_i[\mathcal Q_i(\theta_c)-\theta_c\cdot Z_{c,i}]
       =-\sum_i\partial_i\theta_c\cdot Z_{c,i}.
 \]
 Substitute the strong flux approximation already proved. The remainder
 satisfies
 \[
 \|r_c^{\mathcal E}\|_{L^1_{t,X}}
 \le\|\nabla_X\theta_c\|_{L^2_{t,X}}
     \|Z_c-c^{-1}\mathsf K(N_c)\nabla_X\theta_c\|_{L^2_{t,X}}
       =o(c^{-1}).
 \]
 This proves \eqref{cur:macroentropy}.
 Proposition \ref{ons:tensor} gives nonnegativity of the leading
 dissipation.
\end{proof}

\begin{remark}[Topology, initial layer and normalization]\label{cur:scope}
The time domain of the strong convergence is $(0,T)$, and its norms
cover the entire prescribed interval. At the initial time $z_c(0)=0$,
whereas the solution of \eqref{cur:limitcell} can be nonzero there.
This is why the conclusions use time-integrated topologies.
The momentum norm retains the degenerate weight in $H_\nu$;
the cell solution need not belong to unweighted $L^2(D)$.
The free-transport corner traces of Theorem \ref{hyd:main} are preserved.
Equation \eqref{cur:firstorder} describes the first-order flux of the
actual solution. A limit for $c(\theta_c-\bar\theta)$ requires further
analysis. Under $c=4\pi/\varepsilon$, the tensor coefficient becomes
$\varepsilon\mathsf K/(4\pi)$. No rate for strong convergence of the
resonance current is asserted here.
\end{remark}
